\documentclass[UTF8,fontset=none,zihao=-4,a4paper]{ctexart}
\usepackage[margin=25mm]{geometry}
\setCJKmainfont{Noto Serif CJK SC}[BoldFont=Noto Serif CJK SC Bold]
\setCJKsansfont{Noto Sans CJK SC}
\setmainfont{Times New Roman}
\usepackage[unicode,hidelinks,pdfauthor={},pdftitle={AI工具使用详情}]{hyperref}
\setlength{\parskip}{3pt}
\linespread{1.18}
\begin{document}
\begin{center}{\Large\bfseries AI工具使用详情}\end{center}
\noindent 本说明对应当前论文及支撑材料；不包含参赛者、学校或赛区身份。

\section*{1. 工具名称与型号}
可确认本次使用OpenAI Codex。历史会话的完整工具清单、版本或型号尚未形成可核对记录，需参赛队依据实际使用记录补充；不把当前工具名称推定为全部历史会话的型号。

\section*{2. 使用目的与环节}
根据参赛队说明，AI主要用于选题与研究思路讨论、论文草稿整理、代码初步草拟。核心建模与分析由参赛队主导，模型假设、方案选择、结果解释及最终取舍由参赛队负责；AI提供辅助建议，不替代参赛成员的判断。

实际可见工作还包括既有实验材料整理、正文与图注草拟和修订、绘图与统计复核脚本编写、结果一致性检查、引用和版式核查，以及支撑材料打包。这些环节属于辅助流程，不应缩写为仅做语言润色。正式测试表依据已有日志整理，未生成或替换实验数据。

\section*{3. 主要提示方式与过程}
提示以题目条件、已有模型、代码、实验结果和队员的修改要求为输入，围绕候选思路比较、推导表述、草稿结构、代码框架与错误排查展开；再依据现有证据修订文字、图表和程序。当前提示重点是按format2026整理支撑材料，在参考文献前说明AI用途，并突出人工主导。

论文整理时区分解析证明、有限场景检查、自建实验与官方运行记录；保留失败案例和无显著收益的候选，不用AI生成的数据代替原始记录。源程序、配置、必要结果及正式日志随支撑材料保存。

\section*{4. 输出采纳、人工修改与核验}
参赛队对AI提供的思路、文字及代码初稿进行了人工筛选与修改。选题与方案讨论中，队员结合题目条件确定建模方向、模型假设和求解方法；论文整理中，对草稿的逻辑结构、数学表述及结果解释进行了修订；代码实现中，对初步代码进行了补充、调试，并结合测试结果检查其与模型的一致性。队员依据原始实验记录核对了主要数据、图表及结论，对缺乏依据或不适用的建议未予采纳。核心建模、方案取舍和最终结论由参赛队自主确定，AI仅作为辅助工具，参赛队对论文和程序的最终内容负责。

上述人工处理情况由参赛队确认。自动辅助检查包括统计复算、文件哈希和排版检查，作为人工核验的辅助，不替代参赛成员的判断。官方正式日志保持原名与原始字节，未对加密行为体进行解密或修改。
\end{document}
